\documentclass[12pt,addpoints,answers]{exam}
\usepackage[utf8]{inputenc}
\usepackage{amsmath,amsfonts,amssymb,amsthm}
\usepackage[margin=1in]{geometry}
\usepackage{mathtools}
\usepackage{hyperref}
\usepackage{fullpage}
\usepackage{microtype}
\usepackage{xspace}
\usepackage[svgnames]{xcolor}
\usepackage[sc]{mathpazo}
\usepackage{enumitem}
\usepackage{bm}
\usepackage[mathscr]{euscript}

%----------Header--------------------%
\def\course{{\sc Introduccion a la Criptografia Moderna}}
\def\term{Depto. de Computaci\'{o}n, FCEN, UBA, 1er. Bimestre 2026}
\def\prof{Lecturer: Juan Garay}
\newcommand{\handout}[5]{
   \renewcommand{\thepage}{\arabic{page}}
   \begin{center}
   \framebox{
      \vbox{
    \hbox to 5.78in { \hfill \large{\course} \hfill }
    \vspace{2mm}
%    \hbox to 5.78in { \hfill \large{\prof} \hfill }
%       \vspace{2mm}
       \hbox to 5.78in { {\Large \hfill \textbf{#5}  \hfill} }
       \vspace{2mm}
       \hbox to 5.78in { \term \hfill \emph{#2}}
       \hbox to 5.78in { {#3 \hfill \emph{#4}}}
      }
   }
   \end{center}
   \vspace*{4mm}
}
\newcommand{\hw}[4]{\handout{#1}{#2}{#3}{#4}{Homework #1}}

% -- For ignoring stuff -- %
\newcommand{\ignore}[1]{}


\begin{document}

%----Specs: change accordingly-----%
\newif\ifstudent % comment out false
\studenttrue 
% \studentfalse

\def\hwnum{3}
\def\issuedate{28/4/26}
\def\duedate{12/5/26, 23:59 hs} % 
\def\yourname{Pedro Fuentes Urfeig} % put your name here
%------------------------------%

\ifstudent
\hw{\hwnum}{\issuedate}{Student: \yourname}{Due: \duedate}%
\else
\hw{\hwnum}{\issuedate}{\prof}{Due: \duedate}%
\fi

\noindent \textbf{Instructions}

\begin{itemize}
    \item Upload your solution to Campus; make sure it's only one file, and clearly write your name on the first page. Name the file \textsf{`$<$your last name$>$\_HW3.pdf.'} 
    %{\bf Important:} Make sure to tap {\bf Turn in} after you upload your solution.   
      
     If you are proficient with \LaTeX, you may also typeset your submission and submit in PDF format. To do so, uncomment the ``\%\textbackslash begin\{solution\}'' and ``\%\textbackslash end\{solution\}'' lines and write your solution between those two command lines.
    
      \item Your solutions will be graded on \emph{correctness} and
    \emph{clarity}. You should only submit work that you believe to be
    correct.
    % , and you will get significantly more partial credit if you     clearly identify the gap(s) in your solution.
    
    \item You may collaborate with others on this problem set.  However,
    you must \textbf{{write up your own solutions}} and \textbf{{list
      your collaborators and any external sources (including ChatGPT and similar generative AI chatbots)}} for each
    problem. Be ready to explain your solutions orally to a member of the course's staff
    if asked.
    
    \ignore{
    \item For problems that require you to provide an algorithm, you must
    give a precise description of the algorithm, together with a proof
    of correctness and an analysis of its running time. You may use
    algorithms from class as subroutines. You may also use any facts
    that we proved in class or from the book.
    } %IGNORE
    
\end{itemize}


\noindent This assignment contains \numquestions\ questions,
% \numpages\ pages
for a total of \numpoints \ points.
% and \numbonuspoints\ bonus points.

\bigskip
%\newpage

\begin{questions}

\question[10] Consider the following key-exchange protocol:
\begin{enumerate}[label=(\alph*)]
    \item Alice chooses uniform $k,r\in\{0,1\}^n$, and sends $s:=k\oplus r$ to Bob.
    \item Bob chooses uniform $t\in\{0,1\}^n$, and sends $u:=s\oplus t$ to Alice.
    \item Alice computes $w:=u\oplus r$ and sends $w$ to Bob.
    \item Alice outputs $k$ and Bob outputs $w\oplus t$.
\end{enumerate}
Show that Alice and Bob output the same key. Analyze the security of the scheme (i.e., either prove its security or show a concrete attack).

\begin{solution}
  Veamos que computan la misma clave. 

  $k_A := k$

  $k_B := w \oplus t = u \oplus r \oplus t = s \oplus t \oplus r \oplus t = s \oplus r = k \oplus r \oplus r = k$ 

  Veamos ahora que no es seguro. Si un atacante está viendo los mensajes que se intercambian Alice y Bob, puede hacer lo siguiente. 
  En la primera comunicación, donde Alice le envía $s$, el adversario se lo guarda.
  Luego, Cuando Bob envía $u$, el adversario puede computar $t := s \oplus u$.
  Finalmente, cuando Bob envía $w$, el adversario computa $k_B := w \oplus t$. Y así consigue armar la clave de Bob y, por lo tanto, de Alice. 

\end{solution}

\question[10] Formally define the CDH assumption (i.e., first define a CDH experiment, and then use it to define what it means for the CDH problem to be hard). Prove that hardness of the CDH problem relative to $\mathcal{G}$ implies hardness of the discrete-logarithm problem relative to $\mathcal{G}$, and that hardness of the DDH problem relative to $\mathcal{G}$ implies hardness of the CDH problem relative to $\mathcal{G}$.
\begin{solution}

  Definimos el siguiente experimento $CDH$
  \begin{enumerate}
    \item Usar $\mathcal{G}(1^n)$ para obtener $(\mathbb{G}, q, g)$. 
    \item Elegir $x_1$, $x_2$ de manera uniforme. 
    \item Enviarle al adversario $\mathscr{A}(\mathbb{G}, q, g, g^{x_1}, g^{x_2})$ 
    \item Si el adversario devuelve $g^{x_1 \cdot x_2}$, el experimento termina con resultado 1. Si no, 0. 
  \end{enumerate}

  Definimos CDH de la siguiente manera. El problema CDH es difícil relativo a $\mathcal{G}$ si el experimento descrito arriba tiene probabilidades negligibles de tener éxito. Es decir, 
  $$Pr(CDH(\mathcal{G}, 1^n)) \leq negl(n)$$

  Ahora, demostremos que la dificultad del problema CDH (relativa a $\mathcal{G}$) implica dificultad del problema de logaritmo discreto (relativa a $\mathcal{G}$). Podemos hacer la demostración por contrarrecíproco. Es decir, queremos demostrar que si el logaritmo discreto es fácil, CDH también lo es. 
Asumimos que tenemos un adversario $\mathscr{A}_{DLOG}$ que resuelve el problema del logaritmo discreto con probabilidad no negligible (es decir, dados $(\mathbb{G}, q, g, h)$, puede devolver $x$ tal que $g^x = h$). 
  Luego, definimos un adversario $\mathscr{A}_{CDH}$ tal que hace lo siguiente. 
  \begin{enumerate}
    \item Recibe $(\mathbb{G}, q, g, g^{x_1}, g^{x_2})$. 
    \item Obtiene $x_1 := \mathscr{A}_{DLOG}(\mathbb{G}, q, g, g^{x_1})$
    \item Devuelve $(g^{x_2})^{x_1}$
  \end{enumerate}

  Vemos que este adversario tiene probabilidad no negligible de ganar el experimento, ya que su probabilidad de éxito depende únicamente de la probabilidad (no negligible) de que $\mathscr{A}_{DLOG}$ tenga éxito. 



  Ahora, demostremos que la dificultad del problema DDH (relativo a $\mathcal{G}$) implica la dificultad del problema CDH (relativo a $\mathcal{G}$). 
  Vemos que el problema DDH implica que un adversario distinga entre $g^z$ (con $z$ uniforme) y $g^{x_1 \cdot x_2}$, siendo $x_1$, $x_2$ elegidos de manera uniforme. La definición es:

  $$\left| \Pr[\mathcal{A}(\mathbb{G}, q, g, g^{x_1}, g^{x_2}, g^z) = 1] - \Pr[\mathcal{A}(\mathbb{G}, q, g, g^{x_1}, g^{x_2}, g^{x_1\cdot x_2}) = 1] \right| \leq \mathsf{negl}(n)$$

   Veamos la demostración por el contrarrecíproco. Quiero ver que, si el problema CDH es fácil con respecto a $\mathcal{G}$, entonces el problema DDH es fácil con respecto a $\mathcal{G}$. 

   Asumimos entonces que tenemos un algoritmo que resuelve CDH con probabilidad no negligible, y queremos resolver DDH con probabilidad no negligible. 
   Construimos un adversario $\mathscr{A}_{DDH}$ que hace lo siguiente. 
   \begin{enumerate}
    \item Recibe $(\mathbb{G}, q, g, g^{x_1}, g^{x_2}, g^k)$. El adversario no sabe si $k$ es un valor uniforme cualquiera, o si $k = x_1 \cdot x_2$.
    \item Obtiene $x_1 := CDH(\mathbb{G}, q, g, g^{x_1}$. CDH tiene éxito con probabilidad no negligible. 
    \item Luego, computa $h := (g^{x_2})^{x_1}$, puesto que $g^{x_2}$ es un parámetro que recibe. Luego, si CDH tuvo éxito, $h = g^{x_1 \cdot x_2}$
    \item Finalmente, devuelve $1$ si $g^k = g^{x_1 \cdot x_2}$, $0$ si no.
   \end{enumerate}

   Vemos que la probabilidad de que este adversario tenga éxito es la misma probabilidad de que CDH tenga éxito. Luego, este adversario tiene éxito con probabilidad no negligible de distinguir entre un $g^k$ con $k$ arbitrario y $g^{x_1 \cdot x_2}$

\end{solution}


\question Suppose you know two El Gamal ciphertexts $\langle c_1,c_2\rangle$ and $\langle c'_1,c'_2\rangle$ that encrypt unknown plaintexts $m$ and $m'$. You also know the public key $h$ and cyclic group generator $g$.
\begin{parts}
    \part[5] What information can you infer about $m$ and $m'$ if you observe that $c_1=c'_1$?
    \begin{solution}
      Dada la construcción de los ciphertexts de El Gamal, podemos ver que la única razón para que eso sea cierto (asumiendo que ambos ciphertexts son generados con el mismo $g$) es porque $y = y'$, donde $y$ es un elemento de $\mathbb{Z}_q$ elegido de manera uniforme en la llamada a $ENC$. 

      Luego, veamos qué consecuencias tiene esto para $c_2$ y $c'_2$

      $c'_2 = h^y \cdot m' = \frac{c_2}{m} \cdot m'$. Luego, $\frac{m'}{m} = \frac{c'_2}{c_2}$.

    \end{solution}
    \part[5] What information can you infer about $m$ and $m'$ if you observe that $c_1=g\cdot c'_1$?
    \begin{solution}
      Por construcción de $c_1$ y $c'_1$, sabemos esto.


      $c_1 = g^y = g \cdot c'_1 = g \cdot g^{y'} = g^{y'+1}$


      Luego, $y = y' + 1$. Veamos qué implica esto para $c_2 y c'-2$. 

      $c'_2 = h^{y-1} \cdot m' = h^y \cdot h^{-1} \cdot m' = \frac{c_2}{m} \cdot h^{-1} \cdot m'$

      Si reacomodamos, $\frac{m'}{m} = \frac{c'_2 \cdot h}{c_2}$
    \end{solution}
\end{parts}


\question[10] In class we showed that El Gamal encryption is {\em malleable}, and specifically that given a ciphertext $\langle c_1,c_2\rangle$ that is the encryption of some unknown message $m$, it is possible to produce a ciphertext $\langle c_1,c'_2\rangle$ that is the encryption of $\alpha\cdot m$ (for known $\alpha$). A receiver who receives both these ciphertexts might be suspicious since both ciphertexts share the first component. Show that it is possible to generate $\langle c'_1,c'_2\rangle$ that is the encryption of $\alpha\cdot m$, with $c'_1\neq c_1$ and $c'_2\neq c_2$.
\begin{solution}

\end{solution}


\question[10] Explain why the RSA exponents $e$ and $d$ must always be odd numbers.
% \begin{solution}
% % Uncomment and type your solution here
% \end{solution}


\question[10] Consider the RSA-based encryption scheme in which a user encrypts a message $m\in\{0,1\}^\ell$ with respect to the public key $\langle N,e\rangle$ by computing $\tilde{m}:=H(m)\parallel m$ and outputting the ciphertext $[\tilde{m}^e~\mathrm{mod}~N]$. (Here, let $H:\{0,1\}^\ell\rightarrow\{0,1\}^n$ and assume $\ell+n<\Vert N\Vert$.) The receiver recovers $\tilde{m}$ in the usual way and verifies that it has the correct form before outputting the $\ell$ least-significant bits as $m$. Prove or disprove that this scheme is CCA-secure if $H$ is modeled as a random oracle.
% \begin{solution}
% % Uncomment and type your solution here
% \end{solution}



\end{questions}

\end{document}
